\documentclass{article}
\usepackage{graphicx} % Required for inserting images
\usepackage{amsmath}
\usepackage{amssymb}

\title{iterative-linear-quadratic-regulator}
\author{huangzhengming28@gmail.com}
\date{February 2026}

\begin{document}

\maketitle

\section{Introduction}
This project is based on the basic implementation of linear-quadratic-regulator(LQR). LQR is applcated to linear time-invariant system. For nonlinear system, we can use LQR to find the optimal control law. The main idea is to transform the problem into a "differential world". In this world, the problem is equivalent to LQR. So the main work here is about differential.

\section{Differential}
1. Nonlinear system dynamics
\begin{equation}
    x_{t+1} = f(x_t, u_t)
\end{equation}
\begin{equation}
    \delta{x_{t+1}} = \frac{\partial{f}}{\partial{x_t}} \delta{x_t} + \frac{\partial{f}}{\partial{u_t}} \delta{u_t} = A \delta{x_t} + B \delta{u_t}
\end{equation}
2. Q value of $x_t, u_t$
\begin{gather}
    \begin{bmatrix}
        \delta{x_t}\\
        \delta{u_t}
    \end{bmatrix} = \delta{X}\\
    \delta{Q(x_t, u_t)} = \frac{1}{2} \delta{X}^T \begin{bmatrix}Q_{xx}, Q_{xu}\\Q_{ux}, Q_{uu} \end{bmatrix} \delta{X} + [Q_x, Q_u] \delta{X} \label{q_value_function}
\end{gather}
3. V value of $x_t$
\begin{gather}
    \delta{V(x_{t})} = \delta{Q(x_{t})} + \delta{V(x_{t+1})}\\
    \delta{Q(x_{t}, u_t)} = \delta(x_t^T Q x_t + u_t^T R u_t)=[lx, lu]\delta{X} + \frac{1}{2} \delta{X}^T \begin{bmatrix}
        l_{xx}, l_{xu}\\
        l_{ux}, l_{uu}
    \end{bmatrix} \delta{X}\\
    \delta{V(x_{t+1})} = [p_{t+1}^T A, p_{t+1}^T B] \delta{X} + \frac{1}{2} \delta{X}^T \begin{bmatrix}
        A^T P_{t+1} A, A^T P_{t+1} B\\
        B^T P_{t+1} A, B^T P_{t+1} B
    \end{bmatrix} \delta{X}
\end{gather}
Therefore
\begin{equation}
    \delta{V(x_{t})} = \frac{1}{2} \delta{X}^T \begin{bmatrix}
        l_{xx} + A^T P_{t+1} A, l_{xu} + A^T P_{t+1} B\\
        l_{ux} + B^T P_{t+1} A, l_{uu} + B^T P_{t+1} B
    \end{bmatrix} \delta{X} + [l_x + p_{t+1}^T A, l_u + p_{t+1}^T B] \delta{X} \label{v_value_function}
\end{equation}

\section{Algorithm}
The algorithm of iLQR is the same as LQR. The key values is optimal feedback gain $K_t$, Q-value-coefficient such as $Q_{uu}$ and V-value-coefficient such as $P_{t+1}$.
\subsection{Solve optimal control command}
By focusing on $\delta{u_t}$ with \eqref{q_value_function}, it's optimal value is solved in the same way as LOR:
\begin{gather}
    \delta{\hat{u_t}} = -Q_{uu}^{-1}(Q_{ux}x_t + Q_u) = K_t \delta{x_t} + d_t \label{optimal_control_command}\\
    K_t = -Q_{uu}^{-1} Q_{ux}\\
    d_t = -Q_{uu}^{-1} Q_u
\end{gather}

\subsection{Derive Q-value-coefficients}
With optimal control command $\delta{u_t}$, \eqref{q_value_function} equals to \eqref{v_value_function}, which is:
\begin{equation}
    \frac{1}{2} \delta{X}^T \begin{bmatrix}Q_{xx}, Q_{xu}\\Q_{ux}, Q_{uu} \end{bmatrix} \delta{X} + [Q_x, Q_u] \delta{X}\\
    = \frac{1}{2} \delta{X}^T \begin{bmatrix}
        l_{xx} + A^T P_{t+1} A, l_{xu} + A^T P_{t+1} B\\
        l_{ux} + B^T P_{t+1} A, l_{uu} + B^T P_{t+1} B
    \end{bmatrix} \delta{X} + [l_x + p_{t+1}^T A, l_u + p_{t+1}^T B] \delta{X}
\end{equation}
So the Q-value-coefficients are derived:
\begin{gather}
    Q_x = l_x + p_{t+1}^T A\\
    Q_u = l_u + p_{t+1}^T B\\
    Q_{xx} = l_{xx} + A^T P_{t+1} A\\
    Q_{xu} = l_{xu} + A^T P_{t+1} B\\
    Q_{ux} = l_{ux} + B^T P_{t+1} A\\
    Q_{uu} = l_{uu} + B^T P_{t+1} B
\end{gather}

\subsection{Derive V-value-coefficient}
$p_t, P_t$ are used to represent $V(x_t)$ after $\hat{u_t}$ is solved. They are used for next recursion step. With \eqref{q_value_function} by taking in \eqref{optimal_control_command} and focusing on $x_t$:
\begin{gather}
    \delta{V(x_t)} = \delta{x_t}^T(Q_{xx} + Q_{xu} K_t + K_t^T Q_{ux} + K_t^TQ_{uu}K_t)\delta{x_t} + (Q_x + Q_u K_t + Q_{xu} d_t + K_t O_{uu} d_t) \delta{x_t}\\
    = \delta{x_t}^T P_t \delta{x_t} + p_t \delta{x_t}
\end{gather}
So the V-value-coefficients are derived:
\begin{gather}
    P_t = Q_{xx} + Q_{xu} K_t + K_t^T Q_{ux} + K_t^TQ_{uu}K_t\\
    p_t = Q_x + Q_u K_t + Q_{xu} d_t + K_t O_{uu} d_t
\end{gather}
Final $p_N$ and $P_N$:
\begin{gather}
    p_N = \frac{\partial{x_N^T Q x_N}}{\partial{x_N}}\\
    P_N = \frac{\partial^2{x_N^T Q x_N}}{\partial{x_N}^2}
\end{gather}

\subsection{Forward Pass}
Forward Pass is different from LOR. For each iteration j, the state is $\delta{x_t} = x_t^{j} - x_t^{j-1}$, which is the differential between two iterations. The result is:
\begin{gather}
    \hat{\delta{u_t}} = K_t \delta{x_t} + d_t\\
    u_t^j = u_t^{j-1} + \hat{\delta{u_t}} \\
    x_{t+1}^j = f(x_t^j, u_t^j)
\end{gather}

\subsection{Iteration Convergence Condition}
Given a $\epsilon$ and metric $J^i$, iteration ends when:
\begin{equation}
    |\frac{J^{j} - J^{j-1}}{J^{j-1}}| < \epsilon
\end{equation}

\end{document}
